\documentclass{article}
\usepackage[UTF8]{ctex}
\usepackage{amsmath,amssymb,amsfonts}
\usepackage{enumitem}
\usepackage{geometry}
\geometry{a4paper, margin=1in}
% === 错题本排版优化（不参与题目标识） ===
\setlength{\parindent}{0pt}
\setlength{\parskip}{0.35em plus 0.10em minus 0.08em}
\setlength{\emergencystretch}{3em}
\setlist{itemsep=0.25em, topsep=0.40em, parsep=0pt}
\allowdisplaybreaks[2]
% === 错题本排版优化结束 ===
\begin{document}
\section*{复旦大学2025年数学分析真题}
\section*{34．复旦大学2025年数学分析部分}

\subsection{题目1}
1．求极限 \(\lim _{n \rightarrow \infty}\left(\frac{1+\sqrt[n]{2025}}{2}\right)^{n}=\) \(\_\_\_\_\)．

\subsection{题目2}
2 ．求不定积分 \(\int \frac{1+\sin x}{1+\cos x} e^{x} \mathrm{~d} x=\) \(\_\_\_\_\)．

\begin{itemize}

\subsection{题目3}
\end{itemize}

\subsection{题目3}
原题边界“3．”：\(f(x)\) 为定义在 \((0,1)\) 上的实值函数，满足
\[
\lim _{x \rightarrow 0^{+}} f(x)=0, \lim _{x \rightarrow 0^{+}} \frac{f(x)-f\left(\frac{x}{2}\right)}{x}=0 .
\]
则 \(\lim _{x \rightarrow 0^{+}} \frac{f(x)}{x}=\) \(\_\_\_\_\)．

\subsection{题目4}

\subsection{题目4}
原题边界“4．”：平面 \(x+y+z=0\) 与椭球面 \(x^{2}+y^{2}+4 z^{2}=1\) 相交得到的椭圆面积 \(S=\) \(\_\_\_\_\)．

\subsection{题目5}
5．设 \(L\) 为任一包含原点的光滑曲线，取逆时针方向，求 \(\oint_{L} \frac{x \mathrm{~d} y-y \mathrm{~d} x}{x^{2}+y^{2}}=\) \(\_\_\_\_\)．

二．（20分）设 \(\alpha>\beta>1\) ，求 \(\int_{0}^{+\infty} \frac{1}{1+x^{\alpha}|\sin x|^{\beta}} \mathrm{d} x\) 的敛散性。

三．（20分）求使级数 \(\sum_{n=1}^{\infty}\left(\frac{1}{n} \csc \frac{1}{n}-1\right)^{2025 p}\) 收敛的 \(p\) 的取值范围．
\end{document}
